%!TEX root = ../template.tex
%%%%%%%%%%%%%%%%%%%%%%%%%%%%%%%%%%%%%%%%%%%%%%%%%%%%%%%%%%%%%%%%%%%%
%% chapter2.tex
%% NOVA thesis document file
%%
%% Chapter 2: Background
%%%%%%%%%%%%%%%%%%%%%%%%%%%%%%%%%%%%%%%%%%%%%%%%%%%%%%%%%%%%%%%%%%%%

\typeout{NT FILE chapter2.tex}%

\chapter{Background}
\label{cha:background}

\section{Software modelling}
\label{sec:software_modelling_conceptual}

\subsection{Purpose and role of modelling in software engineering}
\label{subsec:purpose_role_modelling_se}


A software model shows the parts of a system, or of the problem, that matter for a purpose, such as talking about the system or trying a design \cite{kuhne:DagSemProc.04101.15,gorschek2014design,petre2013uml}. The symbols have an agreed meaning: the same mark is read in the same way. The language is the set of allowed symbols and those meanings. A model is one description written in that language for one problem \cite{kuhne:DagSemProc.04101.15}.

\subsection{Conceptual modelling as a broader discipline}
\label{subsec:conceptual_modelling_broader_discipline}

Conceptual modelling is the wider family, including data, process, and goal models. This thesis is about software modelling. The courses use some of those notations, such as UML, BPMN, entity-relationship models, and goal models.

\subsection{Modelling methods, languages, and artefacts}
\label{subsec:modelling_methods_languages_artefacts}

One account of a modelling method includes the language and the steps for creating a valid model from a description. It also includes what a tool can do with that model once it exists, such as simulating it or turning it into another form \cite{bork2019framework}. The language is the allowed symbols and what they mean. What they mean is called the semantics. Students can copy the symbols and still be unable to take those steps or judge the result.

If a model is stored so a machine can read it, it can be queried, analysed, or transformed.

\subsection{Modelling education}
\label{sec:modelling_education}

Modelling education, in this thesis, means the courses in which students learn software modelling: what a construct means, how to build a model from a description, and how to judge whether it is good enough. Modelling is often a smaller part of the programme than programming, and students often treat it as less central than coding courses.

\subsection{Modelling tools}
\label{subsec:tools_mde_education}

A modelling tool is the program a student uses to draw a model, to check it, or to turn it into another form, such as code. In model-driven engineering, the tool treats the model as the source of that other form.

\subsection{Implications for the scope of this thesis}
\label{subsec:implications_scope_thesis}

This thesis is about software modelling in courses, where students and lecturers deal with what a construct means, how to build a model from a description, and how to judge whether it is good enough.

\section{Bloom's revised taxonomy}
\label{sec:blooms_taxonomy}

Bloom's revised taxonomy names six levels: remember, understand, apply, analyse, evaluate, and create \cite{krathwohl2002revision}. This thesis uses four of them for interview prompts: remembering a notation, understanding a construct, applying a procedure, and evaluating whether a model is good enough. Metacognition means noticing one's own way of working.

\section{Grounded Theory and Socio-Technical Grounded Theory}
\label{sec:grounded_theory}

Grounded Theory builds an explanation from the interviews as they are collected \cite{glaser1967discovery,stol2016grounded}. The researcher codes those interviews, writes memos, and compares each new incident with what is already coded.

Socio-Technical Grounded Theory (STGT) uses those same practices in software engineering \cite{hoda2018socio}. The explanation has to cover the people, the practices, and the tools together. A Grounded Theory study can leave the tool as background. An STGT study treats the tool as part of the account.

This study used open coding, memos written during coding, and constant comparison. It did not use the Strauss and Corbin coding paradigm. Interviewing stopped when the evidence was enough for the claims of the study. Sims and Cilliers call that theoretical sufficiency \cite{sims2023qualitatively}.

\section{Human factors}
\label{sec:human_factors}

Human factors are the conditions around a task that change how a person does it. In this thesis they are the conditions around a modelling task in a course: the tool, the subject of the assignment, who reads the model, the feedback, and the load of the course. They are not the notation itself.

Tools shape how easily students can try ideas, get feedback, and recover from mistakes, and a difficult tool can turn students against modelling even when the notation is clear (Section~\ref{subsec:tools_mde_education}).
